% Options for packages loaded elsewhere
\PassOptionsToPackage{unicode}{hyperref}
\PassOptionsToPackage{hyphens}{url}
\PassOptionsToPackage{dvipsnames,svgnames,x11names}{xcolor}
%
\documentclass[
  11pt,
]{article}
\usepackage{amsmath,amssymb}
\usepackage{iftex}
\ifPDFTeX
  \usepackage[T1]{fontenc}
  \usepackage[utf8]{inputenc}
  \usepackage{textcomp} % provide euro and other symbols
\else % if luatex or xetex
  \usepackage{unicode-math} % this also loads fontspec
  \defaultfontfeatures{Scale=MatchLowercase}
  \defaultfontfeatures[\rmfamily]{Ligatures=TeX,Scale=1}
\fi
\usepackage{lmodern}
\ifPDFTeX\else
  % xetex/luatex font selection
\fi
% Use upquote if available, for straight quotes in verbatim environments
\IfFileExists{upquote.sty}{\usepackage{upquote}}{}
\IfFileExists{microtype.sty}{% use microtype if available
  \usepackage[]{microtype}
  \UseMicrotypeSet[protrusion]{basicmath} % disable protrusion for tt fonts
}{}
\makeatletter
\@ifundefined{KOMAClassName}{% if non-KOMA class
  \IfFileExists{parskip.sty}{%
    \usepackage{parskip}
  }{% else
    \setlength{\parindent}{0pt}
    \setlength{\parskip}{6pt plus 2pt minus 1pt}}
}{% if KOMA class
  \KOMAoptions{parskip=half}}
\makeatother
\usepackage{xcolor}
\usepackage[margin=2.4cm]{geometry}
\usepackage{longtable,booktabs,array}
\usepackage{calc} % for calculating minipage widths
% Correct order of tables after \paragraph or \subparagraph
\usepackage{etoolbox}
\makeatletter
\patchcmd\longtable{\par}{\if@noskipsec\mbox{}\fi\par}{}{}
\makeatother
% Allow footnotes in longtable head/foot
\IfFileExists{footnotehyper.sty}{\usepackage{footnotehyper}}{\usepackage{footnote}}
\makesavenoteenv{longtable}
\setlength{\emergencystretch}{3em} % prevent overfull lines
\providecommand{\tightlist}{%
  \setlength{\itemsep}{0pt}\setlength{\parskip}{0pt}}
\setcounter{secnumdepth}{-\maxdimen} % remove section numbering
\usepackage{newunicodechar}
\usepackage{amssymb}
\newunicodechar{±}{\ensuremath{\pm}}
\newunicodechar{²}{\ensuremath{^{2}}}
\newunicodechar{³}{\ensuremath{^{3}}}
\newunicodechar{·}{\ensuremath{\cdot}}
\newunicodechar{¹}{\ensuremath{^{1}}}
\newunicodechar{×}{\ensuremath{\times}}
\newunicodechar{ő}{\H{o}}
\newunicodechar{Σ}{\ensuremath{\Sigma}}
\newunicodechar{α}{\ensuremath{\alpha}}
\newunicodechar{β}{\ensuremath{\beta}}
\newunicodechar{δ}{\ensuremath{\delta}}
\newunicodechar{η}{\ensuremath{\eta}}
\newunicodechar{θ}{\ensuremath{\theta}}
\newunicodechar{ι}{\ensuremath{\iota}}
\newunicodechar{κ}{\ensuremath{\kappa}}
\newunicodechar{λ}{\ensuremath{\lambda}}
\newunicodechar{μ}{\ensuremath{\mu}}
\newunicodechar{ν}{\ensuremath{\nu}}
\newunicodechar{π}{\ensuremath{\pi}}
\newunicodechar{ρ}{\ensuremath{\rho}}
\newunicodechar{σ}{\ensuremath{\sigma}}
\newunicodechar{τ}{\ensuremath{\tau}}
\newunicodechar{ψ}{\ensuremath{\psi}}
\newunicodechar{ω}{\ensuremath{\omega}}
\newunicodechar{…}{\ldots{}}
\newunicodechar{′}{\ensuremath{'}}
\newunicodechar{″}{\ensuremath{''}}
\newunicodechar{⁰}{\ensuremath{^{0}}}
\newunicodechar{⁴}{\ensuremath{^{4}}}
\newunicodechar{⁵}{\ensuremath{^{5}}}
\newunicodechar{⁶}{\ensuremath{^{6}}}
\newunicodechar{⁷}{\ensuremath{^{7}}}
\newunicodechar{⁸}{\ensuremath{^{8}}}
\newunicodechar{⁹}{\ensuremath{^{9}}}
\newunicodechar{⁺}{\ensuremath{^{+}}}
\newunicodechar{⁻}{\ensuremath{^{-}}}
\newunicodechar{₀}{\ensuremath{_{0}}}
\newunicodechar{₁}{\ensuremath{_{1}}}
\newunicodechar{₂}{\ensuremath{_{2}}}
\newunicodechar{₃}{\ensuremath{_{3}}}
\newunicodechar{₄}{\ensuremath{_{4}}}
\newunicodechar{₅}{\ensuremath{_{5}}}
\newunicodechar{₆}{\ensuremath{_{6}}}
\newunicodechar{₇}{\ensuremath{_{7}}}
\newunicodechar{₈}{\ensuremath{_{8}}}
\newunicodechar{ℓ}{\ensuremath{\ell}}
\newunicodechar{ℛ}{\ensuremath{\mathcal{R}}}
\newunicodechar{→}{\ensuremath{\to}}
\newunicodechar{↦}{\ensuremath{\mapsto}}
\newunicodechar{∅}{\ensuremath{\emptyset}}
\newunicodechar{∈}{\ensuremath{\in}}
\newunicodechar{∉}{\ensuremath{\notin}}
\newunicodechar{∎}{\ensuremath{\blacksquare}}
\newunicodechar{∏}{\ensuremath{\prod}}
\newunicodechar{−}{\ensuremath{-}}
\newunicodechar{∖}{\ensuremath{\setminus}}
\newunicodechar{√}{\ensuremath{\surd}}
\newunicodechar{∞}{\ensuremath{\infty}}
\newunicodechar{∤}{\ensuremath{\nmid}}
\newunicodechar{∥}{\ensuremath{\parallel}}
\newunicodechar{∩}{\ensuremath{\cap}}
\newunicodechar{∪}{\ensuremath{\cup}}
\newunicodechar{≈}{\ensuremath{\approx}}
\newunicodechar{≠}{\ensuremath{\neq}}
\newunicodechar{≡}{\ensuremath{\equiv}}
\newunicodechar{≢}{\ensuremath{\not\equiv}}
\newunicodechar{≤}{\ensuremath{\le}}
\newunicodechar{≥}{\ensuremath{\ge}}
\newunicodechar{⊄}{\ensuremath{\not\subseteq}}
\newunicodechar{⊆}{\ensuremath{\subseteq}}
\newunicodechar{⊔}{\ensuremath{\sqcup}}
\newunicodechar{⌊}{\ensuremath{\lfloor}}
\newunicodechar{⌋}{\ensuremath{\rfloor}}
\newunicodechar{𝒞}{\ensuremath{\mathcal{C}}}
\newunicodechar{𝒫}{\ensuremath{\mathcal{P}}}
\newunicodechar{∣}{\ensuremath{\mid}}
\newunicodechar{∗}{\ensuremath{\ast}}
\newunicodechar{⋯}{\ensuremath{\cdots}}
\newunicodechar{≪}{\ensuremath{\ll}}
\newunicodechar{∼}{\ensuremath{\sim}}
\ifLuaTeX
  \usepackage{selnolig}  % disable illegal ligatures
\fi
\IfFileExists{bookmark.sty}{\usepackage{bookmark}}{\usepackage{hyperref}}
\IfFileExists{xurl.sty}{\usepackage{xurl}}{} % add URL line breaks if available
\urlstyle{same}
\hypersetup{
  pdftitle={Stability, rigidity and exact extremisers for Erds Problem 848},
  pdfauthor={Luke Culpitt (AI assisted)},
  colorlinks=true,
  linkcolor={blue},
  filecolor={Maroon},
  citecolor={Blue},
  urlcolor={Blue},
  pdfcreator={LaTeX via pandoc}}

\title{Stability, rigidity and exact extremisers for Erd\H{o}s Problem
848}
\author{Luke Culpitt (AI assisted)}
\date{November 2026}

\begin{document}
\maketitle

\subsection{Abstract}

Erdős and Sárközy asked whether the largest A ⊆ \{1, \ldots, N\} with ab
+ 1 never squarefree is the residue class 7 mod 25. Sawhney proved this
for large N, and recent computer-assisted work, notably Li's
Lean-verified proof, extends it to every N. We give an independent proof
for every N by a different method, with explicit threshold N₀ =
6.7838·10⁶. It rests on a finite ``witness game'' on 63 residue types
modulo 900, whose value 0.0294059\ldots{} is certified by exact integer
branch-and-bound.

Our main new results describe valid sets beyond the maximum:

\begin{itemize}
\tightlist
\item
  \textbf{Explicit stability.} For N ≥ 6.7838·10⁶, every valid set not
  inside one class has at most (N − 7)/25 elements, and at most 0.0293N
  asymptotically. This makes Sawhney's qualitative η explicit: η =
  0.0016 at N = 10⁷, rising to about 0.0107 as N → ∞.
\item
  \textbf{Sharp rigidity.} One element outside 7 mod 25 leaves at most a
  fraction 0.3667\ldots{} of that class. This is the sharp form of the
  degree estimate in Li's proof.
\item
  \textbf{Exact maximisers for every N.} They are the classes 7 and 18
  mod 25, plus four small sets not inside one class when 41 ≤ N ≤ 81.
\item
  \textbf{A two-class set of density 1/50.} It meets both classes, it
  beats the natural ``one class plus one outsider'' configuration, and
  it is optimal. Every valid set not inside one class has at most N/50 +
  O(N\textsuperscript{0.81}) elements, so the second-best density is
  exactly 1/50. Near that density F is rigid: for large N, a valid set
  not inside one class with at least (1/50 − 0.0008)N elements lies
  inside F or its mirror. A companion paper extends the method to ab + 1
  never k-free, to ab + c never squarefree, and to the third density
  (Section 9).
\end{itemize}

The constructions and the certificate for the second-best density are
verified in Lean 4. Every computation is reproducible and was
cross-checked by an independent implementation. This includes a check of
the key degree proposition of Li's proof for 5·10⁶ ≤ N ≤ 10⁷.

\subsection{1. Introduction}

Call A ⊆ {[}N{]} = \{1, \ldots, N\} \textbf{valid} if ab + 1 is not
squarefree for all a, b ∈ A, including a = b. Erdős and Sárközy asked
whether the maximum of \textbar A\textbar{} over valid sets is attained
by \{n ≤ N : n ≡ 7 (mod 25)\}
(\href{https://www.erdosproblems.com/848}{Erdős Problem \#848}). This
class is valid because 7² + 1 = 50, so 25 \textbar{} ab + 1 whenever a ≡
b ≡ 7 (mod 25). The class 18 ≡ −7 works for the same reason.

Van Doorn gave \textbar A\textbar{} ≤ (0.108\ldots{} + o(1))N, and
Weisenberg improved the constant to about 0.105 by inclusion--exclusion.
Early computations posted by StijnC already pointed to stability, and
found small valid sets avoiding both classes, such as \{41, 70, 239,
437\} at N = 606. Sawhney
(\href{https://www.math.columbia.edu/~msawhney/Problem~848~.pdf}{note})
proved the conjecture for N ≥ N₀, with N₀ not computed, and noted a
qualitative stability property. Sawhney and Sellke discuss the problem
in Section IV.1 of arXiv:2511.16072, and Alexeev, Mixon and Sawin treat
the companion problem where ab is never squarefree (arXiv:2507.01928).
As Alexeev pointed out, Erdős's original question has no restriction to
large N. Later AI-assisted work on the forum, led by Sothanaphan, made
N₀ explicit, reaching exp(1420)
(\href{https://www.erdosproblems.com/forum/thread/848}{forum}). That is
far beyond the range of computation.

Since mid-2026 several unrefereed, AI-assisted all-N claims have
appeared. Sothanaphan reached N₀ = 2.64·10¹⁷, and full claims came from
Pitchford with Codex
(\href{https://zenodo.org/records/21647629}{Zenodo}), A. C. Li (SSRN,
with Lean; analytic range reportedly from 5·10⁶), Zanghi
(\href{https://github.com/bzanghi/erdos-848}{GitHub}, N₀ = 5·10⁸) and
Durfee
(\href{https://github.com/AstroCat93/erdos-848-verification}{GitHub}, N₀
= 4.2·10⁷). All of these follow Sawhney's case split, and most cover
small N with colouring certificates. The argument here is independent of
them. Its new ingredients are the type-matrix witness game, which gives
N₀ = 6.78·10⁶, and Theorems B and C, which none of those claims state.

Write c7(N) = \#\{n ≤ N : n ≡ 7 (mod 25)\} = ⌊(N − 7)/25⌋ + 1, and C7,
C18 for the classes 7 and 18 mod 25.

\textbf{Theorem A (Li; independent proof here).} For every N ≥ 1, every
valid A ⊆ {[}N{]} has \textbar A\textbar{} ≤ c7(N).

\textbf{Theorem B (stability).} Let N ≥ 6.7838·10⁶. If A ⊆ {[}N{]} is
valid and not contained in C7 or C18, then \textbar A\textbar{} ≤ (N −
7)/25 \textless{} c7(N). Moreover \textbar A\textbar{} ≤ (0.029289 +
o(1))N. Equivalently, every valid set with more than about 0.0293N
elements lies in one class.

\textbf{Theorem C (rigidity and extremisers).}

\begin{enumerate}
\def\labelenumi{\arabic{enumi}.}
\tightlist
\item
  If A is valid and A ⊄ C7, then \textbar A ∩ C7\textbar{} ≤ (g1* +
  o(1))N, where g1* = (1/3 + 2/3·β∞)/25 = 0.0146697\ldots{} and β∞ = 1 −
  ∏\textsubscript{p≥7}(1 − p⁻²). This constant is sharp.
\item
  For N ∉ {[}41, 81{]}, the maximisers in Theorem A are exactly C7 ∩
  {[}N{]}, together with C18 ∩ {[}N{]} when c18(N) = c7(N). For 41 ≤ N ≤
  81 there are also the maximisers \{7, 41\}, \{32, 43\}, \{41, 43\} and
  \{32, 43, 57\}, each in its own range of N.
\end{enumerate}

\textbf{Theorem D (a dense two-class set).} The set F = \{n ≡ 7 (mod
25), n ≡ 1 (mod 4)\} ∪ \{n ≡ 18 (mod 25), n ≡ 3 (mod 4)\} is valid. It
has density 1/50 and meets both classes. Among sets defined by
congruences mod 900, it is the densest valid set not inside one class.
Moreover F is optimal: every valid A ⊆ {[}N{]} not inside C7 or C18 has
\textbar A\textbar{} ≤ N/50 + 30·N\textsuperscript{0.81}, so the
second-best density λ* equals 1/50 (Section 8). For N ≥
10\textsuperscript{23.6}, every valid set not inside one class with at
least (1/50 − 0.0008)N elements lies inside F or its mirror (Theorem
8.9).

Theorems B, C and D are new. Sawhney's note gives stability with an
unspecified η and leaves open whether 18 mod 25 also attains the
maximum. Li's proof of Theorem A, and the other all-N claims, give no
stability constant and no description of the maximisers. The Lean
formalisation in erdos-banger proves stability with η = 1/2000 beyond an
ineffective N₀.

\textbf{Method.} Each a lies in one of 63 \emph{types}, given by its
class mod 25 (C7, C18 or the rest R), mod 4 and mod 9. Whether p²
\textbar{} ab + 1 is forced for p ∈ \{2, 3, 5\} depends only on the
types of a and b. Fix a valid A and pick one or two ``witnesses'' b ∈ A.
Every other a ∈ A must make ab + 1 non-squarefree. Optimising the
witness choice against every possible set of types gives a finite
min--max problem with value V₂ = 0.0294059\ldots{} \textless{} 1/25.

Sections 3--5 bound the error terms explicitly, giving Theorem A for N ≥
6.7838·10⁶. Section 6 covers smaller N by computation. Section 7 proves
Theorems B and C and the first part of Theorem D, and Section 8 proves
that F is optimal and, near 1/50, essentially unique. Section 9
summarises further results from a companion paper, and Section 10 lists
open problems.

The arguments and code were developed with AI assistance and
independently re-run (Appendix). They have not yet been reviewed by a
human expert.

\subsection{2. Local structure}

Let V = \{x ≥ 1 : x² + 1 is not squarefree\}. Taking a = b shows that
every valid set lies in V.

\textbf{Lemma 2.1.} x ∈ V if and only if p² \textbar{} x² + 1 for some
prime p ≡ 1 (mod 4). For p = 5 this means x ≡ 7 or 18 (mod 25). Hence C7
∪ C18 ⊆ V, and x ∈ R lies in V iff p² \textbar{} x² + 1 for some prime p
≡ 1 (mod 4) with p ≥ 13.

\emph{Proof.} x² + 1 ≡ 1 or 2 (mod 4) and is never divisible by 3. An
odd prime p divides x² + 1 only if −1 is a square mod p, that is, p ≡ 1
(mod 4). The roots of x² ≡ −1 (mod 25) are 7 and 18. ∎

\textbf{Types.} The \emph{type} of a is the triple (κ, μ, ν) where:

\begin{itemize}
\tightlist
\item
  κ ∈ \{C7, C18, R\} is the class of a mod 25;
\item
  μ ∈ \{E, 1, 3\} records a even, a ≡ 1 (mod 4) or a ≡ 3 (mod 4);
\item
  ν ∈ \{0, 1, 2, 4, 5, 7, 8\} records 3 \textbar{} a, or the unit class
  of a mod 9.
\end{itemize}

This gives 3·3·7 = 63 types. Each is a union of residue classes mod 900.
Type t has density w\textsubscript{t} = w(κ)·w(μ)·w(ν), where w(C7) =
w(C18) = 1/25, w(R) = 23/25, w(E) = 1/2, w(1) = w(3) = 1/4, w(0) = 1/3,
and w(ν) = 1/9 otherwise.

\textbf{Forcing.} For q² ∈ \{4, 9, 25\}, say a witness type s
\emph{forces} a target type t if q² \textbar{} ab + 1 for every a of
type t and b of type s, for some such q. This happens exactly when one
of the following holds:

\begin{longtable}[]{@{}ll@{}}
\toprule\noalign{}
Modulus & Forced when \\
\midrule\noalign{}
\endhead
\bottomrule\noalign{}
\endlastfoot
25 & a, b both in C7, or both in C18 \\
4 & a, b both odd, in different classes mod 4 \\
9 & a, b both units mod 9, with ab ≡ 8 (mod 9) \\
\end{longtable}

When t is a C-type and s does not force t, no a of type t and b of type
s has 4 \textbar{} ab + 1, 9 \textbar{} ab + 1 or 25 \textbar{} ab + 1.
For 25, note that a ≡ 7 and 25 \textbar{} ab + 1 force b ≡ 7, and
similarly for 18. So non-squarefreeness must then come from p² with p ≥
7.

\textbf{Lemma 2.2 (forced share).} If b ∉ C7, the C7-types forced by b
have total density at most (1/25)(1 − (3/4)(8/9)) = (1/3)(1/25).

\emph{Proof.} Forcing by 25 is impossible. Forcing by 4 removes at most
one of the four classes mod 4, and forcing by 9 at most one of the nine
classes mod 9. These conditions are independent. ∎

The class 7 mod 25 is special because a single modulus (25) forces every
pair inside it. Every other source of non-squarefreeness has density
well below 1.

\subsection{3. Counting lemmas}

Throughout, U is a union of types and w\textsubscript{U} its density.
Two kinds of event are counted:

\begin{itemize}
\tightlist
\item
  the \emph{B-event} p² \textbar{} ab + 1, for a fixed witness b and a
  prime p ≥ 7 with p ∤ b;
\item
  the \emph{A-event} p² \textbar{} a² + 1, for a prime p ≡ 1 (mod 4)
  with p ≥ 13.
\end{itemize}

The B-event fixes a in one class mod p², namely a ≡ −b⁻¹. The A-event
fixes a in two classes, ±ι\textsubscript{p}.

\textbf{Lemma 3.1 (discrepancy).} Let U = S × \{c\}, with S a set of
residues mod 36 and c a class mod 25. Let q ≥ 1 be coprime to 30. Then
for any residue r,

\[\#\{a \le N : a \in U,\ a \equiv r \ (q)\} = \frac{|S|\,N}{900\,q} + \theta \cdot 19, \qquad |\theta| \le 1.\]

The error is at most 38 for a union over C7 and C18, at most 57 for S ×
R, and at most 3 for all of R. Write e\textsubscript{U} for the bound
that applies to U.

\emph{Proof.} Write a = a₀ + 25q·i. Since 25q is a unit mod 36, the
condition a ∈ S (mod 36) holds for i in a set of \textbar S\textbar{}
classes mod 36. A run of J consecutive i hits it J\textbar S\textbar/36
+ O(18) times, and J ≤ N/(25q) + 1. For R, write it as Z/25 minus the
two C classes. ∎

\textbf{Lemma 3.2 (small primes).} Let 𝒫 be a finite set of primes not
dividing 30b, and let event p take r\textsubscript{p} classes mod p².
Then

\[\#\{a \le N,\ a\in U : \text{some event } p\in\mathcal P \text{ holds}\} = w_U N\Big(1-\prod_{p\in\mathcal P}\big(1-\tfrac{r_p}{p^2}\big)\Big) + \theta\, e_U \Big(1+\prod_{p\in\mathcal P}(1+r_p)\Big).\]

\emph{Proof.} Take the complement and apply Möbius inversion over
subsets of 𝒫. Each term is a union of ∏ r\textsubscript{p} classes, each
counted by Lemma 3.1. ∎

\textbf{Lemma 3.3 (medium primes).} For a single prime p, the count of a
∈ U with the B-event exceeds w\textsubscript{U} N/p² by at most
min(e\textsubscript{U}, 2N/(25p²) + 2, N/p² + 1). It exceeds it by at
most 1 if p² \textgreater{} N. For the A-event on R the excess is at
most min(6, 2(N/p² + 1)), and at most 1 if p² \textgreater{} 2N.

\emph{Proof.} These are Lemma 3.1, then the count ignoring the mod-36
condition, then the count ignoring U. A class mod p² \textgreater{} N
meets {[}N{]} at most once. The roots satisfy ι + (p² − ι) = p², so at
most one of them is ≤ N when p² \textgreater{} 2N. ∎

\textbf{Lemma 3.4 (Pell bound for large primes in x² + 1).} Let M = (N²
+ 1)/Y² and m\textsubscript{j} = ((T\textsubscript{j} +
T\textsubscript{j}⁻¹)/2)², where T\textsubscript{j} = (2N +
1)\textsuperscript{1/(2j−1)}. Then

\[\#\{x \le N,\ x\in R : p^2 \mid x^2+1 \text{ for a prime } p>Y\} \le \sum_{j\ge1} \#\{2\le m\le \min(M,m_j) : m \text{ admissible}\}.\]

Here m is \emph{admissible} if m mod 8 ∈ \{1, 2, 5\}, 25 ∤ m, and m has
no prime factor in \{3, 7, 11, 19, 23\}.

\emph{Proof.} Write x² + 1 = mp². Then m \textless{} M, and m ≥ 2. The
number m is not a square, since X² − Z² = −1 has no solution with X ≥ 1.
Primes q ≡ 3 (mod 4) never divide x² + 1. Also x² + 1 ≡ 1, 2, 5 (mod 8)
and p² ≡ 1 (mod 8), and 25 ∤ x² + 1 for x ∈ R.

The positive solutions of X² − mZ² = −1 are η\textsuperscript{2j−1},
where η is the fundamental unit of Z{[}√m{]}. From mZ² = X² + 1 with Z ≥
1 we get η ≥ √(m−1) + √m. Also η\textsuperscript{2j−1} = x + √(x²+1) ≤
2N + 1. So a j-th solution with x ≤ N needs √(m−1) + √m ≤
T\textsubscript{j}, that is, m ≤ m\textsubscript{j}. The map x ↦ (m, j)
is injective. ∎

\textbf{Lemma 3.5 (lattice bound for large primes in ab + 1).} Fix b,
and let T\textsubscript{b} be the C7 ∪ C18 classes a with 4 ∤ ab + 1, 9
∤ ab + 1 and 25 ∤ ab + 1. Let ω₁(b) be the number of distinct primes ≥ 7
dividing b. Then

\[\#\{a\le N,\ a\in T_b : p^2\mid ab+1,\ p>Y\} \le 2^{\omega_1(b)}\Big[\frac{2\lambda(N+1)}{Y} + \frac{c\sqrt{bN+1}}{b} + \frac{\lambda(bN+1)}{Y^2} + c\Big].\]

The constants (λ, c) depend only on the \{2, 3, 5\}-part of b; in all
cases λ ≤ 0.5334 and c ≤ 18.94.

\emph{Proof.} Write ab + 1 = ms² with s = p \textgreater{} Y. Then m
\textless{} (bN + 1)/Y². Since s is coprime to 30, m mod 1800 lies in a
set S\textsubscript{b} determined by b mod 1800. Also ms² ≡ 1 (mod b),
so s lies in one of ρ\textsubscript{b}(m) classes mod b. Here
ρ\textsubscript{b}(m) ≤ ρ\textsubscript{b₀}(m)·2\textsuperscript{ω₁(b)},
with b₀ the \{2, 3, 5\}-part of b. So the count for a given m is at most
ρ\textsubscript{b}(m)(√((bN+1)/m)/b + 1).

The weight g(m) = ρ\textsubscript{b₀}(m)·1{[}m ∈ S\textsubscript{b}{]}
is 1800-periodic. Computed exactly, its partial sums satisfy
Σ\textsubscript{m≤x} g(m) ≤ λx + c.~Partial summation then gives Σ
g(m)/√m ≤ 2λ√M + c.~∎

\subsection{4. The type-matrix bound}

Fix truncation points z\textsubscript{A}, z\textsubscript{B} and put

\[\beta^{+} = 1-\prod_{7\le p\le z_B}\Big(1-\frac1{p^2}\Big) + \sum_{p>z_B}\frac1{p^2}, \qquad \alpha_s = 1-\prod_{\substack{13\le p\le z_A\\ p\equiv1\,(4)}}\Big(1-\frac2{p^2}\Big), \qquad \tau_A = 2\sum_{\substack{p>z_A\\ p\equiv1\,(4)}}\frac1{p^2}.\]

The tails are evaluated exactly up to 10⁷, and beyond that by the bound
Σ\textsubscript{p\textgreater x} p⁻² ≤ 2.51012/(x ln x), which follows
from Rosser--Schoenfeld.

\textbf{The witness game.} Let P be a set of types, and W ⊆ P with 1 ≤
\textbar W\textbar{} ≤ 2. Define

\[\mathrm{val}(P,W) = \sum_{\substack{t\in P\\ t \text{ a C-type}}} w_t\cdot\begin{cases}1 & \text{every } s\in W \text{ forces } t\\ \beta^{+} & \text{otherwise}\end{cases} \;+\; \sum_{\substack{t\in P\\ t\subseteq R}} w_t\,\alpha_s \;+\; \mathbf 1[P\cap R\neq\varnothing]\,\tfrac{23}{25}\tau_A,\]

\[V_2 = \max_{P}\ \min_{W\subseteq P,\ 1\le|W|\le 2} \mathrm{val}(P,W),\]

with the maximum over type sets P not contained in the C7-types and not
contained in the C18-types.

\textbf{Theorem 4.1.} If A ⊆ {[}N{]} is valid and not contained in C7 or
C18, then \textbar A\textbar{} ≤ V₂·N + E\textsubscript{I}(N), where
E\textsubscript{I} is the explicit error of Section 5.

\emph{Proof.} Let P be the set of types met by A. P is not contained in
one class, because the type determines the class. Choose W attaining the
minimum, and for each s ∈ W pick b\textsubscript{s} ∈ A of type s. Each
type t ∈ P is bounded as follows.

\begin{enumerate}
\def\labelenumi{\arabic{enumi}.}
\tightlist
\item
  \textbf{C-types forced by every witness.} Use \textbar A ∩ t\textbar{}
  ≤ \textbar t ∩ {[}N{]}\textbar{} = w\textsubscript{t} N + O(1), by
  Lemma 3.1.
\item
  \textbf{Other C-types.} Assign each one to a witness s ∈ W that does
  not force it, and let U\textsubscript{s} be the union of the types
  assigned to s. By Section 2, every a ∈ A ∩ U\textsubscript{s} has p²
  \textbar{} ab\textsubscript{s} + 1 for some prime p ≥ 7. Lemma 3.2
  handles p ≤ z\textsubscript{B}, Lemma 3.3 handles z\textsubscript{B}
  \textless{} p ≤ Y and Lemma 3.5 handles p \textgreater{} Y. Together
  they give \textbar A ∩ U\textsubscript{s}\textbar{} ≤
  w\textsubscript{Us}·β⁺·N + E\textsubscript{B}(N).
\item
  \textbf{R-types.} Use A ∩ t ⊆ V ∩ t. By Lemma 2.1 the count is bounded
  by A-events for primes p ≥ 13. Lemma 3.2 on the union of the R-types
  in P handles p ≤ z\textsubscript{A}. Lemmas 3.3 and 3.4, applied over
  all of R, handle p \textgreater{} z\textsubscript{A}.
\end{enumerate}

Summing gives val(P, W)·N plus the error terms. ∎

\textbf{Value of the game.} The game is a 0--1 problem over 2⁶³ type
sets. It was solved by a mixed-integer program (HiGHS; Q. Huangfu and J.
A. J. Hall, Math. Prog. Comp. 10 (2018), 119--142), and independently by
an exact branch-and-bound in integer arithmetic. The inputs were scaled
by 10¹⁵ and rounded up, so the certified value is an upper bound. For
every one of the 16 pairs (z\textsubscript{A}, z\textsubscript{B}) ∈
\{17, 29, 37, 41\} × \{13, 17, 19, 23\}, the worst P is the set of all
63 types, and

\[V_2 = \tfrac{2}{25}\beta^{+} + \tfrac{23}{25}(\alpha_s+\tau_A).\]

At (z\textsubscript{A}, z\textsubscript{B}) = (17, 13) this is V₂ ≤
0.0294059212, certified with 1,339,011 branch-and-bound nodes. The
smallest of the 16 values is V₂(41, 23) ≤ 0.0292884.

The value is easy to see for P = all types. Take one witness of type
(C7, E, 0) and one of type (C18, E, 0). A witness that is even and
divisible by 3 never forces by 4 or 9, so each C-class gets density β⁺,
from the witness in the other class. The content of the certificate is
that no adversarial choice of P does better than this.

Since 1/25 − 0.0294059 = 0.0105941 \textgreater{} 0, Theorem A follows
for every N with E\textsubscript{I}(N) ≤ 0.0105941N − 0.28. Section 5
makes this explicit.

\subsection{5. The explicit threshold}

\textbf{The error term.} Let k\textsubscript{A} and k\textsubscript{B}
be the numbers of primes in the small-prime ranges of Lemma 3.2: primes
p ≡ 1 (mod 4) with 13 ≤ p ≤ z\textsubscript{A}, and primes with 7 ≤ p ≤
z\textsubscript{B}. Then

\[E_I(N) = 2\cdot 38\,(1+2^{k_B}) + 38 + 57\,(1+3^{k_A}) + E_A(N) + 2E_B(N),\]

\[E_A(N)=\min_Y\Big[\sum_{z_A<p\le Y}\text{(Lemma 3.3 excess)} + \text{(Lemma 3.4)}\Big],\qquad E_B(N)=\max_{b\le N}\min_Y\Big[\sum_{z_B<p\le Y}\text{(Lemma 3.3 excess)} + \text{(Lemma 3.5)}\Big].\]

The maximum over b runs over the 2-, 3- and 5-profile of b and over k =
ω₁(b). Here ω₁(b) = k forces b ≥ b₀·7·11···p\textsubscript{k}, which
tames the factor 2\textsuperscript{k}.

\textbf{The ω-refinement.} Near N₀ the factor 2\textsuperscript{ω₁(b)}
is the bottleneck, so witnesses with many prime factors are avoided. Let
K\textsubscript{k}(N) = \#\{b ≤ N : b ∈ V, ω₁(b) ≥ k\}, computed exactly
(K₅(10⁷) = 672). Split A = A′ ⊔ A″, where A″ is the elements with ω₁ ≥
k₀.

\begin{itemize}
\tightlist
\item
  If A′ is not inside one class, apply Theorem 4.1 to A′ with witnesses
  from A′. Then \textbar A\textbar{} ≤ V₂N + E\textsubscript{I}(N; ω₁
  \textless{} k₀) + K\textsubscript{k₀}(N).
\item
  Otherwise, say A′ ⊆ C7 and A ⊄ C7. Use one witness b ∈ A″
  \textbackslash{} C7. Lemma 2.2 gives \textbar A ∩ C7\textbar{} ≤ g₁N +
  E\textsubscript{II}(N), with g₁ = (1/3 + (2/3)β⁺)/25 ≈ 0.01469, and
  \textbar A \textbackslash{} C7\textbar{} ≤ K\textsubscript{k₀}(N).
\end{itemize}

The second case has much smaller density, so it never binds.

\textbf{Grid certification.} For fixed parameters (z\textsubscript{A},
z\textsubscript{B}, k₀, Y), every term of E is nondecreasing in N. So it
suffices to cover {[}N₀, 10⁶⁰{]} by cells {[}N\textsubscript{lo},
N\textsubscript{hi}{]} with ratio 1.01 below 10¹³ and 1.1 above, and to
check for each cell that

\[E(N_{hi}) \le \big(\tfrac1{25}-V_2\big)N_{lo} - 0.28 .\]

This gives \textbar A\textbar{} ≤ N/25 − 0.28 = (N − 7)/25 \textless{}
c7(N) throughout the cell. Prime counts up to 4·10⁹ are exact, and
beyond that come from Rosser--Schoenfeld and Dusart bounds.

Two independent implementations agree. The referee grid uses only
(z\textsubscript{A}, z\textsubscript{B}) = (17, 13), and its smallest
slack is +329.8, in the first cell {[}6.7838·10⁶, 6.85·10⁶{]}. The
original grid covers {[}10⁷, 10⁶⁰{]}. It optimises the parameters per
cell and has 2525 cells with no failures. Its slack, as a fraction of N,
rises from 0.0016 at 10⁷ to 0.0106 at 10⁶⁰, and tends to 1/25 −
0.0292884 = 0.0107.

\textbf{Proposition 5.1 (tail).} For N ≥ 10⁶⁰, Theorem 4.1 holds with
E\textsubscript{I}(N) ≤ 1.02·10⁻⁸·N.

\emph{Proof.} Take (z\textsubscript{A}, z\textsubscript{B}) = (17, 13),
Y = N\textsuperscript{2/3} and no ω-refinement. Let L = ln N. The terms
are bounded as follows:

\begin{itemize}
\tightlist
\item
  the inclusion--exclusion constant is 1292;
\item
  Lemma 3.3 on the A side is at most 6π(√(2N)) +
  π(N\textsuperscript{2/3});
\item
  Lemma 3.4 is at most (1 + 10⁻⁴)N\textsuperscript{2/3}·J, where J = (1
  + ln(2N+1)/ln(1+√2))/2 = O(L);
\item
  Lemma 3.3 on the B side is at most 38π(√N) +
  π(N\textsuperscript{2/3});
\item
  Lemma 3.5 is at most 2\textsuperscript{k}·(1.067N\textsuperscript{2/3}
  + 20√N).
\end{itemize}

Robin's bound ω(n) ≤ 1.38402·ln n/ln ln n gives 2\textsuperscript{k} ≤
N\textsuperscript{0.95933/ln~L} ≤ N\textsuperscript{0.1947}. The prime
counts use π(x) \textless{} 1.25506x/ln x. Each term divided by N
decreases for N ≥ 10⁶⁰, and the total at N = 10⁶⁰ is 1.02·10⁻⁸·N. ∎

\textbf{Corollary 5.2.} Theorem A holds for all N ≥ 6.7838·10⁶, and
every valid set not inside one class has \textbar A\textbar{} ≤ (N −
7)/25. (Sawhney suggested on the forum that large-sieve methods should
bring N₀ down to about 10⁶. The witness game reaches the same order by a
different route.)

\subsection{6. Small N}

\textbf{N ≤ 3806: exact clique search.} Valid sets are exactly the
cliques in the graph on V ∩ {[}N{]} with edges \{a, b\} where ab + 1 is
not squarefree. At N = 3806 this graph has 397 vertices and 36,248
edges. The clique number is nondecreasing in N, and c7 is constant on
each block (6 + 25k, 31 + 25k{]}. So it suffices to compute the clique
number at N ≡ 6 (mod 25). It equals c7(N) at every such N ≤ 3806.

\textbf{3806 ≤ N ≤ 10⁷: case bounds.} Let H\textsubscript{t}(b, N) =
\#\{a ∈ V ∩ C\textsubscript{t}, a ≤ N : ab + 1 not squarefree\} for t ∈
\{7, 18\}, and let H\textsubscript{odd}(b, N) count the odd such a in C7
∪ C18. Let v\textsubscript{R}(N) = \textbar V ∩ R ∩ {[}N{]}\textbar, and
v\textsubscript{R}\textsuperscript{odd}(N) its odd part. Every valid A
not inside C7 or C18 falls into exactly one of four cases.

\begin{longtable}[]{@{}
  >{\raggedright\arraybackslash}p{(\columnwidth - 4\tabcolsep) * \real{0.3333}}
  >{\raggedright\arraybackslash}p{(\columnwidth - 4\tabcolsep) * \real{0.3333}}
  >{\raggedright\arraybackslash}p{(\columnwidth - 4\tabcolsep) * \real{0.3333}}@{}}
\toprule\noalign{}
\begin{minipage}[b]{\linewidth}\raggedright
Case
\end{minipage} & \begin{minipage}[b]{\linewidth}\raggedright
Condition on A
\end{minipage} & \begin{minipage}[b]{\linewidth}\raggedright
Bound on \textbar A\textbar{}
\end{minipage} \\
\midrule\noalign{}
\endhead
\bottomrule\noalign{}
\endlastfoot
S1 & A ∩ R contains an even b & v\textsubscript{R} + max over even b ∈ V
∩ R of {[}H₇(b) + H₁₈(b){]} \\
S2 & A ∩ R ≠ ∅ and all odd; A ∩ (C7 ∪ C18) all odd &
v\textsubscript{R}\textsuperscript{odd} + max over odd b ∈ V ∩ R of
H\textsubscript{odd}(b) \\
S3 & A ∩ R ≠ ∅ and all odd; an even b′ ∈ C7 ∪ C18 &
v\textsubscript{R}\textsuperscript{odd} + max of H₇(b) over odd b ∈ R,
plus max of H₁₈(b′) over even b′ ∈ C7 (or the same with 7 and 18
swapped) \\
S4 & A ∩ R = ∅, A meets C7 and C18 & max of H₇(b′) over b′ ∈ C18, plus
max of H₁₈(b) over b ∈ C7 \\
\end{longtable}

Each bound is correct because every element of A ∩ C\textsubscript{t} is
compatible with the chosen witness. In S2 an odd R-witness and odd
C-elements remain. In S3 an even C-element of one class witnesses the
other class.

All four bounds are nondecreasing in N. So S(N) = max S\textsubscript{i}
at a block end E\textsubscript{k} bounds every N ∈
(E\textsubscript{k−1}, E\textsubscript{k}{]}. The check is
S(E\textsubscript{k}) ≤ c7(E\textsubscript{k−1} + 1).

With blocks of relative width 1\%, this holds on all 794 blocks from
3806 to 10⁷. For example, S(10⁶) = 37,637 \textless{} c7 = 40,000, and
S(10⁷) = 376,171 \textless{} 400,000. The counts H\textsubscript{t}(b,
N) are computed exactly:

\begin{itemize}
\tightlist
\item
  primes p ≤ (bN)\textsuperscript{1/3} by stepping through the class a ≡
  −b⁻¹ (mod p²);
\item
  larger primes by enumerating ab + 1 = mp² over m, solving s² ≡ m⁻¹
  (mod b) by CRT.
\end{itemize}

Together with Corollary 5.2, whose range starts at 6.7838·10⁶
\textless{} 10⁷, this proves Theorem A for every N.

\subsection{7. Stability, rigidity and the extremisers}

\textbf{Proof of Theorem B.} The first bound is Corollary 5.2. For the
asymptotic bound, apply Theorem 4.1 with (z\textsubscript{A},
z\textsubscript{B}) = (41, 23), where V₂ ≤ 0.0292884. The argument of
Proposition 5.1 gives E\textsubscript{I}(N) = o(N) for any fixed
parameters. ∎

Theorem B implies a robust form. If N ≥ 6.7838·10⁶ and
\textbar A\textbar{} \textgreater{} (N − 7)/25, then A ⊆ C7 or A ⊆ C18,
so A is one full class with some elements removed. Nothing outside the
class can be added. Sawhney's note gives this with an unspecified η.
Here η can be taken to be 0.0016 at N = 10⁷, rising to about 0.0107 as N
→ ∞.

\textbf{Proof of Theorem C(1), upper bound.} Pick b ∈ A \textbackslash{}
C7. Then A ∩ C7 is contained in \{a ∈ C7 : ab + 1 not squarefree\}. By
Lemma 2.2 the forced share s satisfies s ≤ 1/3. On the remaining types
only primes p ≥ 7 contribute, and Lemmas 3.2, 3.3 and 3.5 bound them by
(1 − s)β⁺N/25 plus an error. The bound s + (1 − s)β⁺ increases with s,
so

\[|A\cap C_7| \le \tfrac1{25}\big(\tfrac13+\tfrac23\beta^{+}\big)N + 38\,(1+2^{k_B}) + 38 + E_B(N).\]

This holds for every N. Letting z\textsubscript{B} → ∞ gives the limsup
bound g1*. ∎

\textbf{Sharpness.} Let q be a prime with q ≡ 18 (mod 25), q ≡ 1 (mod 4)
and q ≡ 1 (mod 9). Put A\textsubscript{q}(N) = \{q\} ∪ \{a ∈ C7 ∩
{[}N{]} : aq + 1 not squarefree\}. This set is valid:

\begin{itemize}
\tightlist
\item
  pairs inside C7 have 25 \textbar{} ab + 1;
\item
  the pair (q, q) has 25 \textbar{} q² + 1;
\item
  the pairs (a, q) are valid by definition.
\end{itemize}

The forced share is exactly 1/3, from a ≡ 3 (mod 4) or a ≡ 8 (mod 9). On
the other cells, the events p² \textbar{} aq + 1 for p ≥ 7, p ≠ q are
independent of the cell by CRT. The squarefree sieve for a linear form
gives density β(q) = 1 − ∏\textsubscript{p≥7,~p≠q}(1 − p⁻²). So
\textbar A\textsubscript{q}(N)\textbar/N → (1/3 + (2/3)β(q))/25, which
tends to g1* as q → ∞. ∎

\textbf{Relation to Li's proof.} Li's ``uniform mixed degree''
proposition says that for N ≥ 5·10⁶, every b ∈ V outside C7 ∪ C18 has
more than N/50 + N/525 ≈ 0.0219N squarefree neighbours in C7 ∩ {[}N{]},
and likewise in C18 ∩ {[}N{]}. Theorem C(1) is the sharp form of this.
The squarefree degree is c7(N) − H₇(b), so the worst case is
asymptotically (1/25 − g1*)N ≈ 0.0253N, and no constant above 1/25 − g1*
is possible. Computing H₇ and H₁₈ exactly for every b ≤ 10⁷ confirms
Li's proposition on all of 5·10⁶ ≤ N ≤ 10⁷, with a margin of at least
13.8\%. Each 1\% block is checked conservatively: the largest H at the
top of the block, against the class size at the bottom and Li's
threshold at the top. There max H₇(b)/N stays at 0.01468, matching g1*.

\textbf{Proof of Theorem C(2).} A valid set inside C7 has at most c7(N)
elements, with equality only for C7 ∩ {[}N{]}. The same holds for C18.
So it suffices to compare c7(N) with m(N), the largest size of a valid
set not inside one class. The function m is nondecreasing and c7 is
constant on blocks.

\begin{enumerate}
\def\labelenumi{\arabic{enumi}.}
\tightlist
\item
  \textbf{N ≤ 1006.} Compute m(N) exactly as a maximum over cliques that
  contain an R-vertex or meet both C7 and C18. Then m(N) \textless{}
  c7(N) except for 41 ≤ N ≤ 81, where m = c7. Enumerating all maximal
  cliques for 7 ≤ N ≤ 106 gives exactly the sets in the table below.
\item
  \textbf{1006 \textless{} N ≤ 3806.} The case bounds S1--S4 are
  strictly below c7 on every block.
\item
  \textbf{3806 \textless{} N ≤ 6.9·10⁶.} The case bounds are strictly
  below c7 on all 757 blocks. The smallest block margin is 1, on (4282,
  4324{]}. Checked one N at a time, the margin there is at least 2.
\item
  \textbf{N ≥ 6.7838·10⁶.} This is Theorem B.
\end{enumerate}

\begin{longtable}[]{@{}lll@{}}
\toprule\noalign{}
N & c7(N) & Extra maximisers \\
\midrule\noalign{}
\endhead
\bottomrule\noalign{}
\endlastfoot
41--42 & 2 & \{7, 41\} \\
43--56 & 2 & \{7, 41\}, \{32, 43\}, \{41, 43\} \\
57--81 & 3 & \{32, 43, 57\} \\
\end{longtable}

For instance, 32·43 + 1 = 3⁴·17, 32·57 + 1 = 5²·73 and 43·57 + 1 =
2²·613, and 25 divides 32² + 1, 43² + 1 and 57² + 1. ∎

\textbf{Proof of Theorem D.} If a, b ∈ F lie in the same class mod 25,
then 25 \textbar{} ab + 1. If a ≡ 7 and b ≡ 18 (mod 25), then a ≡ 1 and
b ≡ 3 (mod 4), so ab + 1 ≡ 0 (mod 4). The case a = b is covered by the
first case. F meets both classes, for example 57 ∈ C7 and 43 ∈ C18, and
has density 2·(1/25)(1/4) = 1/50. Maximality among congruence-defined
sets comes from an exact maximum-clique computation on the forcing graph
mod 900. Its value is 1/50, attained by F up to symmetry. The upper
bound for λ* is Theorem 8.1. Validity of F, of C7 and C18, and of the
four exceptional maximisers is checked in Lean 4
(\texttt{Erdos848Constructions.lean}). ∎

\subsection{8. The second-best density}

Let λ*(N) be the largest size of a valid A ⊆ {[}N{]} contained neither
in C7 nor in C18.

\textbf{Theorem 8.1.} For every N, λ*(N) ≤ N/50 + E(N), where E(N) =
N\textsuperscript{2/3+o(1)} is the explicit function of Section 8.5.
Moreover, by the refinement of Section 8.6, λ*(N) ≤ N/50 +
30·N\textsuperscript{0.81} for every N. The set F of Theorem D gives
λ*(N) ≥ N/50 − 2. Hence λ*(N)/N → 1/50.

Section 8.1 proves the theorem for sets inside C7 ∪ C18 by a short
argument. Sections 8.2--8.4 prove the general case: a counting lemma,
then an all-witness version of the game of Section 4, then an exact
certificate. Section 8.5 makes every constant explicit, Section 8.6
lowers the constant to 30, and Section 8.7 proves the third-order
Theorem 8.9.

\subsubsection{8.1 Sets inside C7 ∪ C18}

\textbf{Proposition 8.2.} Let A ⊆ {[}N{]} be valid, A ⊆ C7 ∪ C18, and
suppose A meets both classes. Then \textbar A\textbar{} ≤ N/50 +
O(N\textsuperscript{0.81}).

\emph{Proof.} Put X = A ∩ C7 and Y = A ∩ C18, and let x, y be their
sizes relative to N/25. For a ∈ X and b ∈ Y we have ab + 1 ≡ 7·18 + 1 ≡
2 (mod 25). So ab + 1 is non-squarefree only if:

\begin{itemize}
\item
  \begin{enumerate}
  \def\labelenumi{(\arabic{enumi})}
  \setcounter{enumi}{3}
  \tightlist
  \item
    a and b are odd with ab ≡ 3 (mod 4), or
  \end{enumerate}
\item
  \begin{enumerate}
  \def\labelenumi{(\Alph{enumi})}
  \setcounter{enumi}{16}
  \tightlist
  \item
    9 \textbar{} ab + 1, or p² \textbar{} ab + 1 for some prime p ≥ 7.
  \end{enumerate}
\end{itemize}

For fixed b, let Comp\textsubscript{Q}(b) be the set of a ∈ C7 ∩ {[}N{]}
satisfying (Q). The condition 9 \textbar{} ab + 1 picks out one class
mod 9, and p² \textbar{} ab + 1 one class mod p². Lemma 8.3 handles p
\textgreater{} N\textsuperscript{0.6}. So Comp\textsubscript{Q}(b) has
relative size at most

\[c = \frac19 + \sum_{p\ge 7}\frac1{p^2} = 0.16225\ldots < \frac16,\]

up to O(N\textsuperscript{0.8+o(1)}), and by CRT it meets each class mod
4 proportionally. Split X and Y by residue mod 4 into even, 1 and 3
parts. A pair not forced by 4 lies in Comp\textsubscript{Q} of its
partner. A case analysis on which odd classes each side meets gives:

\begin{itemize}
\tightlist
\item
  one side has no odd element: x + y ≤ 2c;
\item
  one side meets both odd classes: x + y ≤ max(2c, 1/4 + 3c/2);
\item
  both sides meet the same single odd class: x + y ≤ 3c/2;
\item
  they meet opposite single odd classes (the shape of F): x + y ≤ 1/2 if
  neither side has an even element, and otherwise x + y ≤ max(1/4 +
  3c/2, 2c).
\end{itemize}

Since c \textless{} 1/6, every case gives x + y ≤ 1/2. ∎

\subsubsection{8.2 Counting with an arbitrary target}

\textbf{Lemma 8.3 (large primes, any target).} For 1 ≤ b ≤ N and Y ≥ 1,

\[\#\{a\le N : q^2\mid ab+1 \text{ for some prime } q>Y\} \le 2^{\omega(b)+1}\Big(\frac{2(N+1)}{Y}+\frac{N^2+1}{Y^2}\Big).\]

\emph{Proof.} Write ab + 1 = mq². Then 1 ≤ m \textless{} (bN + 1)/Y²,
and ms² ≡ 1 (mod b) with s = q. So s is a square root of m⁻¹ mod b.
There are at most 2\textsuperscript{ω(b)+1} of these: two for each odd
prime power, at most four for a power of 2. Also s ≤ √((bN + 1)/m). So
each m contributes at most 2\textsuperscript{ω(b)+1}(√((bN + 1)/m)/b +
1). Summing over m \textless{} M = (bN + 1)/Y², and using Σ
m\textsuperscript{−1/2} ≤ 2√M, gives 2\textsuperscript{ω(b)+1}(2(bN +
1)/(bY) + M). This is at most the bound stated, since 1 ≤ b ≤ N. ∎

This is a cruder form of Lemma 3.5 that needs no information about the
target. Since 2\textsuperscript{ω(b)} = N\textsuperscript{o(1)}, with Y
= N\textsuperscript{0.6} it is N\textsuperscript{0.8+o(1)}.

\textbf{Lemma 8.4 (joint estimate for R-targets).} Let t be a residue
class mod 900 inside R, of density w\textsubscript{t} (for a type made
of r\textsubscript{t} such classes, add the bounds; each additive error
term is then multiplied by r\textsubscript{t}, as in Section 8.5), and
let b ≤ N be a witness that does not force t. Then

\[\#\{a\in t\cap[N] : a\in V,\ ab+1 \text{ not squarefree}\} \le w_t\, j^{+} N + O(N^{0.8+o(1)}),\qquad j^{+} = \alpha^{+}\beta_{500} + \kappa + \tau + \tau_B \le 0.01577,\]

where

\begin{itemize}
\tightlist
\item
  α⁺ = Σ 2/p² over primes p ≡ 1 (mod 4), p ≥ 13;
\item
  β⁺ = Σ\textsubscript{q≥7} 1/q² = 0.051136 (in this section; it is at
  least the truncated β⁺ of Section 4), and β\textsubscript{500} =
  Σ\textsubscript{7≤q≤500} 1/q² = 0.050863;
\item
  κ = Σ 1/p² over primes p ≡ 1 (mod 4), p ≥ 13;
\item
  τ = 2Σ 1/p² over primes p ≡ 1 (mod 4), p \textgreater{} 500, which is
  at most 0.00028; and τ\textsubscript{B} = Σ 1/q² over all primes q
  \textgreater{} 500, which is also at most 0.00028.
\end{itemize}

\emph{Proof.} Membership a ∈ V needs p² \textbar{} a² + 1 for some p ≡ 1
(mod 4), p ≥ 13 (Lemma 2.1). Because b does not force t,
non-squarefreeness of ab + 1 needs q² \textbar{} ab + 1 for some q ≥ 7.
Split into three parts.

\begin{enumerate}
\def\labelenumi{\arabic{enumi}.}
\tightlist
\item
  \textbf{p, q ≤ 500.} By CRT with mod 900, a pair p ≠ q fixes 2 classes
  mod 900p²q², contributing at most 2(w\textsubscript{t} N/(p²q²) + 1).
  A pair p = q needs a ≡ ±ι\textsubscript{p} and a ≡ −b⁻¹ (mod p²), at
  most one class, contributing w\textsubscript{t} N/p² + 1. Together
  these give at most w\textsubscript{t}(α⁺β\textsubscript{500} + κ)N +
  O(1).
\item
  \textbf{p \textgreater{} 500.} Lemma 3.3 for 500 \textless{} p ≤
  N\textsuperscript{0.6} and Lemma 3.4 for p \textgreater{}
  N\textsuperscript{0.6} give at most w\textsubscript{t} τ N +
  O(N\textsuperscript{0.8} log N).
\item
  \textbf{q \textgreater{} 500.} Drop the condition on p.~For 500
  \textless{} q ≤ N\textsuperscript{0.6} each prime fixes one class,
  contributing at most w\textsubscript{t} N/q² + 1, so these give at
  most w\textsubscript{t} τ\textsubscript{B} N +
  π(N\textsuperscript{0.6}). For q \textgreater{} N\textsuperscript{0.6}
  use Lemma 8.3. ∎
\end{enumerate}

For a C-type target that b does not force, the same argument without the
condition a ∈ V gives at most w\textsubscript{t} β⁺ N +
O(N\textsuperscript{0.8+o(1)}), with β⁺ ≤ 0.05114.

\subsubsection{8.3 The all-witness game and its core}

Use 525 types t: a mod 25, a mod 4 (even, 1 or 3), and a mod 9 with 3
\textbar{} a merged. Let P be the set of types met by a valid A not
inside one class, and choose a witness b\textsubscript{s} ∈ A of each
type s ∈ P. Every a ∈ A ∩ t is compatible with every witness, so
\textbar A ∩ t\textbar{} ≤ min over s ∈ P of the bound for witness
b\textsubscript{s}. The possible bounds are:

\begin{itemize}
\tightlist
\item
  \textbf{Forced C-target:} w\textsubscript{t} N + O(1).
\item
  \textbf{Non-forced C-target:} w\textsubscript{t} β⁺ N +
  O(N\textsuperscript{0.8+o(1)}).
\item
  \textbf{R-target:} w\textsubscript{t} j⁺ N +
  O(N\textsuperscript{0.8+o(1)}). An R-type never forces itself: a² ≡ −1
  has no solution mod 4 or mod 9, and mod 25 only in C7 or C18. So each
  R-type in P can use its own witness.
\end{itemize}

\textbf{Lemma 8.5 (core reduction).} Let K be the set of types in P
forced by every type in P, and let F(K) be the set of types forcing
every member of K. Then K consists of C-types, P ⊆ F(K), and

\[|A| \le \Big(w(K) + \beta\, w\big(\mathcal C\cap F(K)\setminus K\big) + j\, w\big(\mathcal R\cap F(K)\big)\Big)N + O(N^{0.8+o(1)}) =: U(K)\,N + O(N^{0.8+o(1)}),\]

for any β ≥ β⁺ and j ≥ j⁺. Here 𝒞 and ℛ are the C-types and the R-types.

\emph{Proof.} Forcing is symmetric, so every s ∈ P forces every member
of K, giving P ⊆ F(K). A non-core C-type in P is not forced by some s ∈
P, and an R-type in P uses its own witness. Sum the bounds above over t
∈ P. ∎

\subsubsection{8.4 The certificate}

A C-type is forced by a type outside its own class only through 4 or 9,
never 25, since 7·18 + 1 ≡ 2 (mod 25) and 7r + 1 ≢ 0 for r ∈ R. So
forcing between a core and the rest depends only on cells mod 36: 21
cells per class.

A core is a pair (M₇, M₁₈) of cell sets with every cross pair forced.
Each half is automatically internally forced, by 25. U(K) is then a
function of three sets of cells:

\begin{itemize}
\tightlist
\item
  the cells forced with all of M₁₈ (the C7 part of F(K));
\item
  the cells forced with all of M₇ (the C18 part);
\item
  the cells forced with all of M₇ ∪ M₁₈ (the R part, in each of the 23
  classes of R).
\end{itemize}

\textbf{Proposition 8.6.} With β = 0.052 and j = 0.0169, U(K) ≤ 1/50 for
every core K for which F(K) contains a type outside C7 and a type
outside C18. Equality holds only for the core of F and its mirror image.

\emph{Proof.} Write A(M) for the set of cells forced with every cell of
M. The program \texttt{lamcert.c} enumerates all 51,223 such cores and
evaluates 900·U in exact integer arithmetic.

A much smaller check suffices. Replacing M by its closure A(A(M)) does
not change A(M), because A(A(A(M))) = A(M), and it does not lower U,
because the closure contains M. So both halves of the core may be taken
closed. Only 64 sets of cells have the form A(M), and they give 439
admissible pairs.

This reduced check is formalised in \texttt{LamCert.lean} and checked by
the Lean kernel (\texttt{decide}, with no \texttt{native\_decide}). The
file also proves the closure argument, the equality case, and that the
forcing relation on cells is exactly divisibility of ab + 1 by 4 or 9.

The check fails with β = 0.10 and j = 0.03, so it is not vacuous. The
bound survives up to j = 0.01721 (with β = 0.052) or β = 0.0556 (with j
= 0.0169). A floating-point MILP over all 2⁵²⁵ type sets, run
independently, gives the same value 1/50 with a matching dual bound. ∎

\textbf{Proof of Theorem 8.1.} Since β⁺ ≤ 0.05114 \textless{} 0.052 and
j⁺ ≤ 0.01577 \textless{} 0.0169, Lemma 8.5 and Proposition 8.6 give
\textbar A\textbar{} ≤ N/50 + O(N\textsuperscript{0.8+o(1)}), taking Y =
N\textsuperscript{0.6}. Section 8.5 chooses Y better and makes every
term explicit. ∎

\subsubsection{8.5 The explicit error term}

Let r\textsubscript{t} be the number of residues mod 900 in the type t.
These sum to 72 over the 42 C-types, and to 828 over the 483 R-types.
Fix a cut-off Y for the primes q in ab + 1, and Y′ for the primes p in
a² + 1. Going through the proofs of Lemmas 8.3--8.5, the error is at
most the sum of these terms:

\begin{itemize}
\tightlist
\item
  \textbf{Core C-types:} \textbar t ∩ {[}N{]}\textbar{} ≤
  w\textsubscript{t} N + r\textsubscript{t}, so 72 in total.
\item
  \textbf{Other C-types, 7 ≤ q ≤ Y:} one class mod q² for each prime, so
  at most 72π(Y).
\item
  \textbf{Other C-types, q \textgreater{} Y:} a non-forced C-target lies
  in T\textsubscript{b}, so Lemma 3.5 applies. This gives
  42·L\textsubscript{C}, where L\textsubscript{C} is the bound of Lemma
  3.5 with b ≤ N.
\item
  \textbf{R-types, p, q ≤ 500:} 828(2n₁n₂ + n₁), where n₁ = 43 and n₂ =
  92 count the primes in the two ranges.
\item
  \textbf{R-types, 500 \textless{} q ≤ Y:} at most 828·π(Y), one residue
  error per prime q (the condition on p is dropped, as in Lemma 8.4).
\item
  \textbf{R-types, q \textgreater{} Y:} 483·L, where L =
  2\textsuperscript{ω*(N)+1}(2(N+1)/Y + (N²+1)/Y²) comes from Lemma 8.3,
  and ω*(N) is the largest number of distinct prime factors of an
  integer ≤ N.
\item
  \textbf{R-types, 500 \textless{} p ≤ Y′:} at most 1656·π(Y′).
\item
  \textbf{R-types, p \textgreater{} Y′:} the bound of Lemma 3.4, counted
  once for all of R, because the types are disjoint.
\end{itemize}

Let E(N) be the minimum of this sum over Y and Y′. The script
\texttt{explicitC.py} evaluates E on cells of ratio
10\textsuperscript{0.1} up to N = 10⁶⁰¹, refined to ratio
10\textsuperscript{0.01} between 10¹⁰ and 10²⁵. It uses exact prime
counts below 10⁷ and π(x) \textless{} 1.25506x/ln x above. For fixed Y
and Y′ every term is nondecreasing in N, so each cell is checked at its
upper end. Some values:

\begin{longtable}[]{@{}lllllll@{}}
\toprule\noalign{}
N & 10⁷ & 10¹³ & 10²⁰ & 10³⁰ & 10⁶⁰ & 10¹⁰⁰ \\
\midrule\noalign{}
\endhead
\bottomrule\noalign{}
\endlastfoot
E(N)/N & 27 & 0.31 & 0.0030 & 2.4·10⁻⁶ & 4.1·10⁻¹⁵ & 1.1·10⁻²⁶ \\
\end{longtable}

The best Y lies between about N\textsuperscript{0.72} and
N\textsuperscript{0.85}, and E(N) = N\textsuperscript{2/3+o(1)}. The
bound beats the trivial bound \textbar A\textbar{} ≤ N from about N =
10\textsuperscript{11.7}.

So \textbar A\textbar{} ≤ N/50 + C·N\textsuperscript{0.81} holds for
every N, with C the largest value of min((49/50)N\textsuperscript{0.19},
E(N)/N\textsuperscript{0.81}). The grid gives C ≤ 167.5, attained near N
= 10\textsuperscript{11.75}. Beyond 10⁶⁰¹, take Y = Y′ =
N\textsuperscript{2/3} and use Robin's bound ω(n) ≤ 1.38402 ln n/ln ln
n.~This gives E(N) ≤
1001·2\textsuperscript{ω*(N)}N\textsuperscript{2/3}, so
E(N)/N\textsuperscript{0.81} is at most 4·10⁻⁴ and decreasing there.
This gives C = 170. Section 8.6 lowers it to 30.

The constant is large because the witness may have many prime factors,
and because the large-prime count is crude. The exponent 0.81 can be
replaced by any number above 2/3, at the cost of a larger C.

\subsubsection{8.6 Witnesses with many prime factors}

The factor 2\textsuperscript{ω(b)} in Lemmas 3.5 and 8.3 is paid for the
worst possible witness. As in the ω-refinement of Section 5, it can be
avoided by setting aside the elements with many prime factors. Fix k ≥
1, let A″ be the elements of A with at least k distinct prime factors,
and let A′ = A \textbackslash{} A″.

\textbf{Lemma 8.7 (Rankin).} For every z ≥ 1, \#\{n ≤ N : ω(n) ≥ k\} ≤
z\textsuperscript{−k}·N·exp((z − 1)S(N)), where S(N) =
Σ\textsubscript{p≤N} 1/p \textless{} log log N + B + 1/log² N and B =
0.26150\ldots{} is the Meissel--Mertens constant.

\emph{Proof.} Since z\textsuperscript{ω(n)} =
Σ\textsubscript{d~\textbar~n,~d~squarefree} (z −
1)\textsuperscript{ω(d)}, we get Σ\textsubscript{n≤N}
z\textsuperscript{ω(n)} ≤ N·∏\textsubscript{p≤N}(1 + (z − 1)/p) ≤
N·exp((z − 1)S(N)). Divide by z\textsuperscript{k}. The bound on S(N) is
Rosser and Schoenfeld's. ∎

For N ≤ 10¹⁸ a sharper bound comes from Hardy and Ramanujan's induction,
run numerically. Write N\textsubscript{j}(x) = \#\{n ≤ x : ω(n) = j\}
and P(y) for the number of prime powers up to y.

\textbf{Lemma 8.8.} For x ≥ 1 and k ≥ 1,

\[N_{k+1}(x) \le \sum_{\substack{m\le\sqrt{x}\\ \omega(m)=k}} P(x/m) + \frac{1}{k+1}\sum_{\substack{q\le\sqrt{x}\\ q \text{ a prime power}}} N_k(x/q).\]

\emph{Proof.} Let n ≤ x have k + 1 prime-power factors. If the largest
of them, q, exceeds √x, then n = qm with m \textless{} √x and ω(m) = k,
and this representation is unique. Otherwise each of the k + 1
prime-power factors q is at most √x, and n/q has k prime factors and is
at most x/q, so n is counted k + 1 times in the second sum. ∎

The script \texttt{omega\_envelope.py} (in its final form inside
\texttt{explicitC\_omega3.py}) starts from exact counts of
N\textsubscript{j} at the points 10\textsuperscript{i/200} up to 10¹¹.
These come from a segmented sieve, \texttt{segomega.c}, which reproduces
the direct counts of \texttt{omgrid.c} up to 10⁹ exactly. Above that it
uses P(y) ≤ π(y) + 2√y with Dusart's bound for π. It evaluates Lemma 8.8
on the same grid up to 10¹⁸, binning both sums and rounding every
argument up to the next grid point. Since each N\textsubscript{j} is
nondecreasing, this gives upper envelopes. As a check, the same
computation started from exact counts at 10⁷ stays above the exact
counts at every grid point up to 10⁹.

Write K\textsubscript{k}(N) for the smaller of the Rankin and envelope
bounds on \#\{n ≤ N : ω(n) ≥ k\}, with K\textsubscript{k}(N) = 0 when k
\textgreater{} ω*(N). There are two cases.

\begin{itemize}
\tightlist
\item
  \textbf{Case I: A′ is inside neither C7 nor C18.} Run the argument of
  Sections 8.2--8.4 on A′, choosing every witness in A′. Then ω(b) ≤ k −
  1, so in Section 8.5 the factors 2\textsuperscript{ω*(N)+1} and
  2\textsuperscript{ω₁*(N)} become 2\textsuperscript{k} and
  2\textsuperscript{k−1}. Call the resulting error
  E\textsubscript{k}(N). Then \textbar A\textbar{} ≤ N/50 +
  E\textsubscript{k}(N) + K\textsubscript{k}(N).
\item
  \textbf{Case II: A′ ⊆ C7, say.} Since A ⊄ C7, some b ∈ A″ lies outside
  C7, and A \textbackslash{} C7 ⊆ A″. A single cell outside C7 forces at
  most 12 of the 36 residues of C7 mod 36: one row of weight 9 and one
  column of weight 4, which overlap in one cell. So, as in Lemma 2.2,
  \textbar A ∩ C7\textbar{} ≤ g₁N + E\textsubscript{II}(N). Here g₁ =
  (1/3 + (2/3)β⁺)/25 = 0.014697 \textless{} 1/50, and
  E\textsubscript{II}(N) = 36 + 36π(Y) + L\textsubscript{C}, where
  L\textsubscript{C} is Lemma 3.5's bound with the worst-case factor
  2\textsuperscript{ω₁*(N)}. Then \textbar A\textbar{} ≤ N/50 +
  E\textsubscript{II}(N) + K\textsubscript{k}(N) − (1/50 − g₁)N.
\end{itemize}

Two further facts help at moderate N, both because every element of A
lies in V.

\begin{itemize}
\tightlist
\item
  \textbf{Counting inside V.} Let V\textsubscript{P} be the part of V
  that comes from primes p ≤ P = 10⁵, that is, the n with p² \textbar{}
  n² + 1 for some such p.~The sieve also counts the n ∈
  V\textsubscript{P} with ω(n) ≥ k exactly, up to 10¹¹. The rest of V
  comes from primes p \textgreater{} P with p ≡ 1 (mod 4), and has at
  most 2NΣ\textsubscript{p\textgreater P} p⁻² + 2π(Y′) elements up to
  Y′, plus the bound of Lemma 3.4 beyond Y′. Inside V the count of
  elements with many prime factors is about fourteen times smaller than
  over all integers.
\item
  \textbf{The bound \textbar A\textbar{} ≤ \textbar V ∩
  {[}N{]}\textbar.} This replaces the trivial bound \textbar A\textbar{}
  ≤ N. Up to 10¹¹ the sieve gives \textbar V ∩ {[}N{]}\textbar{}
  directly. Beyond that, inclusion--exclusion over the eleven primes p ≤
  100 with p = 5 or p ≡ 1 (mod 4) gives at most 0.10359N + 3¹¹; Lemma
  3.3 covers 100 \textless{} p ≤ Y′, and Lemma 3.4 the primes beyond. At
  10¹¹ this formula gives 1.0549·10¹⁰, against an exact count of
  1.0531·10¹⁰.
\end{itemize}

Taking k = ω*(N) + 1 recovers Section 8.5. The script
\texttt{explicitC\_omega4.py} uses the grid of Section 8.5, refined to
ratio 10\textsuperscript{0.01} from 10⁵ to 10²⁵. For each cell it takes
the better of two bounds. One is \textbar A\textbar{} ≤ \textbar V ∩
{[}N{]}\textbar. The other, for the best k, is the two cases above,
using the smallest available bound on the elements with ω ≥ k. It gives
C ≤ 29.9, attained near N = 10\textsuperscript{13.33} with k = 8. Hence
C = 30 works in Theorem 8.1. For comparison, Rankin's bound alone gives
C ≤ 92.4 (\texttt{explicitC\_omega.py}).

The bound \textbar A\textbar{} ≤ \textbar V ∩ {[}N{]}\textbar{} decides
C up to about 10¹³. Beyond that the refined argument decides it; where C
is decided, about 10¹³ to 10¹⁸, k is between 8 and 11. The worst cell
lies just past the range where the bound \textbar A\textbar{} ≤
\textbar V ∩ {[}N{]}\textbar{} decides C. So extending the sieve, or an
explicit bound for integers with many prime factors in the classes
±ι\textsubscript{p} mod p², would lower C further.

\subsubsection{8.7 Third-order structure}

Write F′ = \{n ≡ 7 (mod 25), n ≡ 3 (mod 4)\} ∪ \{n ≡ 18 (mod 25), n ≡ 1
(mod 4)\} for the mirror image of F.

\textbf{Theorem 8.9.} Let N ≥ 10\textsuperscript{23.6}. If A ⊆ {[}N{]}
is valid, not contained in C7 or C18, and \textbar A\textbar{} ≥ (1/50 −
0.0008)N, then A ⊆ F or A ⊆ F′. The same holds with 0.0005 in place of
0.0008 for N ≥ 10\textsuperscript{22.3}.

So near the second-best density there is no room to move. A
near-extremal set is F or F′ with some elements removed; nothing can be
added.

\textbf{Lemma 8.10 (strict gap).} With β = 0.05114 and j = 0.01577,
every admissible core K other than those of F and F′ has U(K) ≤ 1/50 −
0.0010542. The runner-up is the core of F with one cell of weight 1
removed, which costs exactly (1 − β)/900.

\emph{Proof.} \texttt{lamgap.c}, a variant of \texttt{lamcert.c}, finds
this by enumerating all 51,223 admissible cores in exact integer
arithmetic. \texttt{LamGap.lean} proves the statement in the Lean
kernel, again with \texttt{decide} only. The argument is the closure
argument of Proposition 8.6. If a core's closure is not the core of F or
F′, the 64 × 64 check bounds it. If the closure is F's but the core is
smaller, the core misses a cell of weight at least 1. ∎

\emph{Proof of Theorem 8.9.} The constants of Lemma 8.10 are at least β⁺
and j⁺ (Lemma 8.4). So Lemma 8.5 gives \textbar A\textbar{} ≤ U(K)N +
E(N), where K is the core of A and E(N) is the error of Section 8.5. The
script \texttt{thirdorder.py} checks on the grid of Section 8.5 that
E(N) \textless{} (0.0010542 − δ)N for every N ≥ N₂(δ). It finds N₂ =
10\textsuperscript{23.56} for δ = 0.0008 and N₂ =
10\textsuperscript{22.27} for δ = 0.0005, with the tail beyond 10⁶⁰¹
handled as in Section 8.5. Then U(K) \textgreater{} 1/50 − 0.0010542, so
K is the core of F or of F′; say F.

Every type of A is forced by every type of this core. A C7 cell forced
with every C18 cell ≡ 3 (mod 4) must be ≡ 1 (mod 4), because the cell
with 3 \textbar{} a forces nothing through 9. Symmetrically, a C18 cell
must be ≡ 3 (mod 4). An R-type would have to be forced with the cells ≡
1 and ≡ 3 (mod 4) that have 3 \textbar{} a, through 4 alone, which no
residue mod 4 allows. So every type of A is a type of F, and A ⊆ F. ∎

The threshold is large for two reasons: the gap is small, and the error
of Section 8.5 decays only like N\textsuperscript{−1/3+o(1)}. The
ω-refinement of Section 8.6 does not transfer directly. It would only
show that A lies in F apart from at most K\textsubscript{k}(N)
exceptional elements.

\subsection{9. Further results}

A companion paper {[}14{]} tests how far the method reaches. It proves
counting lemmas for ab + c with any squarefree c, and uses them in three
directions:

\begin{itemize}
\tightlist
\item
  \textbf{Higher powers.} For ab + 1 never k-free, the maximum is a
  class mod 5\textsuperscript{k} for every k ≥ 3 and large N, with an
  explicit threshold. For 3 ≤ k ≤ 100 this holds for every N, with two
  small exceptions; for cubes the maximisers are listed for every N.
\item
  \textbf{Other forms.} For ab + c never squarefree, the maximum for
  large N is a root class of x² + c ≡ 0 (mod p²), for 49 squarefree c ≤
  100.
\item
  \textbf{The third density.} The densest valid set not inside C7, C18,
  F or F′ has density between 1/25 − 1/(4π²) = 0.01467 and 0.01491, so
  for large N, Theorem 8.9 holds with 0.0192 replaced by 0.01491.
\end{itemize}

\subsection{10. Open problems}

\begin{enumerate}
\def\labelenumi{\arabic{enumi}.}
\tightlist
\item
  \textbf{Third-order structure.} Theorem 8.9 settles this for N ≥
  10\textsuperscript{23.6}: a valid set not inside one class with at
  least (1/50 − 0.0008)N elements lies inside F or its mirror. Does this
  hold for much smaller N, with an explicit δ? The exact maximisers for
  N ≤ 1006 (Section 7, proof of Theorem C(2)) are small sets unrelated
  to F, so some lower bound on N is needed.
\item
  \textbf{Sharper large-prime counts.} Brute force shows that Lemmas 3.4
  and 3.5 overestimate by a factor of about 38 or more (worst ratio
  0.026). Averaging ρ\textsubscript{b}(m) over short ranges of m, using
  explicit Burgess- or Weil-type bounds, should lower N₀ substantially,
  make the ω-refinement unnecessary, and shrink the constant C of
  Theorem 8.1.
\item
  \textbf{A general witness-game theorem.} The companion paper {[}14{]}
  carries the method to higher powers and to ab + c for many c, each by
  a separate computation. Is there a single theorem: a self-forcing
  class plus a positive local gap implies extremality for all large N,
  with an effective threshold? The companion paper also asks for the
  exact third density, conjectured to be 1/25 − 1/(4π²).
\item
  \textbf{Formalisation.} The constructions
  (\texttt{Erdos848Constructions.lean}), the certificate of Proposition
  8.6 (\texttt{LamCert.lean}) and the strict gap of Lemma 8.10
  (\texttt{LamGap.lean}) are verified in Lean. Natural next targets are
  Proposition 8.2, which is short, and the counting lemmas of Section 3.
\end{enumerate}

\subsection{Appendix: verification and reproducibility}

All code is plain C (gcc -O3, optionally -fopenmp) or Python 3 with
numpy, sympy, networkx and scipy. Each computational claim for Theorems
A--D has at least two independent checks.

\begin{longtable}[]{@{}
  >{\raggedright\arraybackslash}p{(\columnwidth - 4\tabcolsep) * \real{0.3333}}
  >{\raggedright\arraybackslash}p{(\columnwidth - 4\tabcolsep) * \real{0.3333}}
  >{\raggedright\arraybackslash}p{(\columnwidth - 4\tabcolsep) * \real{0.3333}}@{}}
\toprule\noalign{}
\begin{minipage}[b]{\linewidth}\raggedright
Claim
\end{minipage} & \begin{minipage}[b]{\linewidth}\raggedright
Primary computation
\end{minipage} & \begin{minipage}[b]{\linewidth}\raggedright
Independent check
\end{minipage} \\
\midrule\noalign{}
\endhead
\bottomrule\noalign{}
\endlastfoot
Theorem A, N ≤ 3806 & max clique at N ≡ 6 (mod 25), \texttt{e848b.py}
(776 s) & \texttt{nsc.py}: exact cliques for N ≤ 1006, plus case bounds
S1--S4 strictly below c7 for 1006 \textless{} N ≤ 3806
(\texttt{fast848.c} started at N = 31) \\
Theorem A, 3806 ≤ N ≤ 10⁷ & \texttt{fast848.c}, 794 blocks (47 min) &
\texttt{e848cases.c} (different algorithm) to 10⁶; numpy brute force
\texttt{alledges.py}; \texttt{sbrute.py} (direct d² sieve) matches
S1--S4 at all 168 edges ≤ 20,000 \\
V₂ values & MILP, \texttt{vk2.py} (HiGHS) & exact integer
branch-and-bound, \texttt{bb.c}, all 16 parameter pairs \\
Lemma 3.5 constants λ, c & \texttt{mconst2.py} & \texttt{mc.py}
(brute-force ρ) \\
Lemmas 3.4 and 3.5 & proofs above & brute force at N ≤ 3·10⁵,
\texttt{brute.py} / \texttt{brute2.py}: no violations, worst ratio
actual/bound 0.026 \\
Grid, N ≥ 6.7838·10⁶ & \texttt{tmgrid.py}, 2525 cells on {[}10⁷,
10⁶⁰{]}; the referee grid alone covers {[}6.7838·10⁶, 10⁷{]} & referee
\texttt{cellsbig.py} / \texttt{cellsbig2.py}, (17, 13) only, to 10⁶⁰ \\
ω-counts K\textsubscript{k}(N) & \texttt{omegaV.c} & \texttt{kom.c} \\
Proposition 5.1 & closed form above & \texttt{tail60.py} \\
Theorem C(2), N ≤ 1006 & \texttt{nsc.py} (exact cliques by case) &
maximal-clique enumeration for N ≤ 106 \\
Theorem C(2), 3806 \textless{} N ≤ 6.9·10⁶ & \texttt{fast848.c}, strict
comparison (29 min, 2 cores) & \texttt{sbrute.py} on the tightest block,
one N at a time \\
\end{longtable}

Further independent checks:

\begin{longtable}[]{@{}
  >{\raggedright\arraybackslash}p{(\columnwidth - 4\tabcolsep) * \real{0.3333}}
  >{\raggedright\arraybackslash}p{(\columnwidth - 4\tabcolsep) * \real{0.3333}}
  >{\raggedright\arraybackslash}p{(\columnwidth - 4\tabcolsep) * \real{0.3333}}@{}}
\toprule\noalign{}
\begin{minipage}[b]{\linewidth}\raggedright
Claim
\end{minipage} & \begin{minipage}[b]{\linewidth}\raggedright
Computation
\end{minipage} & \begin{minipage}[b]{\linewidth}\raggedright
Result
\end{minipage} \\
\midrule\noalign{}
\endhead
\bottomrule\noalign{}
\endlastfoot
Li's uniform-degree proposition, 5·10⁶ ≤ N ≤ 10⁷ & \texttt{fastli.c}
(exact H₇, H₁₈ for all b; 70 blocks of 1\%, 56 min) & holds on every
block; minimum margin 13.8\% (threshold taken at the top of each block);
brute-force cross-check at N = 20,000 \\
Theorem D and the constructions & \texttt{Erdos848Constructions.lean},
Lean 4.36.0-rc2 with Mathlib & compiles; no \texttt{sorry}; only the
standard axioms \\
Densest congruence-defined set not inside one class & \texttt{wgame4.py}
(exact cliques mod 900) & 1/50, attained by F; exact integer certificate
of the full game: lamcert.c (51,223 cores); kernel-checked reduced
certificate: LamCert.lean (64 closed sets, 439 pairs); third-order gap:
lamgap.c, LamGap.lean, thirdorder.py; explicit C = 30:
explicitC\textsubscript{omega4}.py, with exact counts to 10¹¹ from
segomega.c (C = 92.4 with Rankin's bound alone:
explicitC\textsubscript{omega}.py; C = 170 without the ω-refinement:
explicitC.py) \\
\end{longtable}

The full package (code, logs and certificates) accompanies this paper.

\subsection{References}

\begin{enumerate}
\def\labelenumi{\arabic{enumi}.}
\tightlist
\item
  T. F. Bloom, \href{https://www.erdosproblems.com/848}{Erdős Problem
  \#848}, erdosproblems.com, and the
  \href{https://www.erdosproblems.com/forum/thread/848}{discussion
  thread}.
\item
  M. Sawhney,
  \href{https://www.math.columbia.edu/~msawhney/Problem~848~.pdf}{Problem
  848}, note.
\item
  W. van Doorn, upper bound (0.108\ldots{} + o(1))N, as recorded at
  reference 1.
\item
  J. B. Rosser and L. Schoenfeld, Approximate formulas for some
  functions of prime numbers, Illinois J. Math. 6 (1962), 64--94.
\item
  G. Robin, Estimation de la fonction de Tchebychef θ sur le k-ième
  nombre premier et grandes valeurs de la fonction ω(n) nombre de
  diviseurs premiers de n, Acta Arith. 42 (1983), no. 4, 367--389,
  doi:10.4064/aa-42-4-367-389.
\item
  P. Dusart, Inégalités explicites pour ψ(X), θ(X), π(X) et les nombres
  premiers, C. R. Math. Rep.~Acad. Sci. Canada 21 (1999), no. 2, 53--59.
  This is the source of π(x) ≤ (x/ln x)(1 + 1/ln x + 2.51/ln² x) for x ≥
  355991.
\item
  B. Alexeev, D. G. Mixon and W. Sawin,
  \href{https://arxiv.org/abs/2507.01928}{The independence and clique
  cover numbers of the squarefree graph} (2025), for the companion
  problem where ab is never squarefree.
\end{enumerate}

\begin{enumerate}
\def\labelenumi{\arabic{enumi})}
\setcounter{enumi}{7}
\tightlist
\item
  A. C. Li, \emph{The Exact Extremal Bound in Erdős Problem 848}, August
  2026
  (\href{https://papers.ssrn.com/sol3/papers.cfm?abstract~id~=7230480}{SSRN
  7230480};
  \href{https://github.com/crabsatellite/erdos-848-squarefree-product}{code
  and Lean}). Uses a Hall-matching reformulation relative to C7, a
  prefix colouring for N ≤ 5·10⁶, and case-wise pivot estimates beyond
  that. Claims a complete Lean 4 verification.
\item
  N. Sothanaphan, explicit threshold N₀ = 2.64·10¹⁷ for Sawhney's
  theorem (2026), as cited in reference 8.
\item
  I. Pitchford (with OpenAI Codex),
  \href{https://zenodo.org/records/21647629}{Candidate
  certificate-backed all-N determination for Erdős Problem 848}, Zenodo,
  July 2026.
\item
  B. Zanghi, \href{https://github.com/bzanghi/erdos-848}{erdos-848}, and
  N. Durfee,
  \href{https://github.com/AstroCat93/erdos-848-verification}{erdos-848-verification},
  GitHub, 2026.
\end{enumerate}

\begin{enumerate}
\def\labelenumi{\arabic{enumi}.}
\setcounter{enumi}{11}
\tightlist
\item
  R. Jung et al.,
  \href{https://github.com/The-Obstacle-Is-The-Way/erdos-banger}{erdos-banger},
  Lean 4 formalisation of Sawhney's stability theorem for Problem 848 (η
  = 1/2000, ineffective N₀), January 2026.
\item
  M. Sawhney and M. Sellke, in \emph{Early science acceleration
  experiments with GPT-5},
  \href{https://arxiv.org/abs/2511.16072}{arXiv:2511.16072}, Section
  IV.1 (Problem 848).
\end{enumerate}

\begin{enumerate}
\def\labelenumi{\arabic{enumi})}
\setcounter{enumi}{13}
\tightlist
\item
  L. Culpitt, \emph{The witness game beyond Erdős Problem 848: higher
  powers, ab + c, and the third density}, companion paper (2026); code
  in the repository
  \href{https://github.com/MeaCulpitt/Erdos-848}{Erdos-848}.
\end{enumerate}

\subsection{Acknowledgements}

The argument, code and certificates were developed with AI assistance
(Claude) and checked by independent re-implementation. Every computation
above can be re-run from the accompanying package. We thank the
contributors to the forum thread for Problem 848, especially W. van
Doorn, D. Weisenberg, StijnC, B. Alexeev, M. Sawhney and N. Sothanaphan,
whose comments shaped this work. We welcome expert review, especially of
Sections 3--5 and 8.

\end{document}
